\begin{derivation}[从\eqnrefs{eq:far_field_source_amplitude}到\eqnrefs{eq:far_field_radiated_energy_green_si}]
真空远场的 Fourier 分量满足

\begin{equation*}
 \widetilde{\mathbf H}_{\rm far}(\omega)
 =\frac1{\mu_0c}\hat{\mathbf n}\times
 \widetilde{\mathbf E}_{\rm far}(\omega),
 \qquad
 \hat{\mathbf n}\cdot\widetilde{\mathbf E}_{\rm far}=0.
\end{equation*}
时域总辐射能为
\begin{align*}
 \frac{\dd W_{\rm rad}}{\dd\Omega}
 &=r^2\int_{-\infty}^{\infty}\dd t\,
 \hat{\mathbf n}\cdot
 [\mathbf E_{\rm far}(t)\times\mathbf H_{\rm far}(t)]\\
 &=\frac{r^2}{\mu_0c}
 \int_{-\infty}^{\infty}\dd t\,|\mathbf E_{\rm far}(t)|^2.
\end{align*}
采用
\begin{equation*}
 \mathbf E(t)=\int\frac{\dd\omega}{2\pi}\,
 \widetilde{\mathbf E}(\omega)\ee^{-\ii\omega t},
\end{equation*}
Parseval 恒等式给
\begin{equation*}
 \int\dd t\,|\mathbf E(t)|^2
 =\int\frac{\dd\omega}{2\pi}|\widetilde{\mathbf E}(\omega)|^2.
\end{equation*}
实场满足 $\widetilde{\mathbf E}(-\omega)=\widetilde{\mathbf E}^{*}(\omega)$，因此正频部分为
\begin{equation*}
 \int\dd t\,|\mathbf E(t)|^2
 =\frac1\pi\int_0^\infty\dd\omega\,
 |\widetilde{\mathbf E}(\omega)|^2.
\end{equation*}
正文的远场源--场关系是
\begin{equation*}
 \widetilde{\mathbf E}_{\rm far}
 =\ii\mu_0\omega\frac{\ee^{\ii\kappa_0r}}{r}
 \mathbf A_\infty,
 \qquad \kappa_0=\omega/c.
\end{equation*}
代入 Poynting 谱后，$r^{-2}$ 与球面面积因子相消，得到
\begin{equation*}
 \frac{\dd^2 W_{\rm rad}}{\dd\omega\dd\Omega}
 =\frac{\mu_0\omega^2}{\pi c}|\mathbf A_\infty|^2,
\end{equation*}
即正文\eqnrefs{eq:far_field_radiated_energy_green_si}。正文\eqnrefs{eq:cl_photon_probability_far_green} 随后只是把这一能量谱除以单光子能量 $\hbar\omega$ 的定义转换，不需要另行引入动力学近似。
\end{derivation}

\begin{derivation}[从\eqnrefs{eq:cl_photon_probability_far_green}到\eqnrefs{eq:free_space_current_radiation_probability}]
真空标量 Helmholtz Green 函数
$g_0(\mathbf r-\mathbf r')=\ee^{\ii\kappa_0|\mathbf r-\mathbf r'|}/[4\pi|\mathbf r-\mathbf r'|]$ 的远场系数由严格极限定义：
\begin{equation*}
\lim_{r\to\infty}
4\pi r\,\ee^{-\ii\kappa_0r}
 g_0(r\hat{\mathbf n}-\mathbf r')
=\ee^{-\ii\kappa_0\hat{\mathbf n}\cdot\mathbf r'}.
\end{equation*}
电磁 Green 张量在同一极限只保留横向投影，因此
\[
\mathbf F_{\infty,0}
(\hat{\mathbf n},\mathbf r';\omega)
=\frac1{4\pi}
(\mathbf I-\hat{\mathbf n}\hat{\mathbf n})
\ee^{-\ii\kappa_0\hat{\mathbf n}\cdot\mathbf r'}.
\]
代入 $\mathbf A_\infty$ 定义：
\begin{align*}
\mathbf A_{\infty,0}
&=\int\dd^3 r'\,
\mathbf F_{\infty,0}\cdot\widetilde{\mathbf J}(\mathbf r',\omega)\\
&=\frac1{4\pi}
(\mathbf I-\hat{\mathbf n}\hat{\mathbf n})
\int\dd^3 r'\,\widetilde{\mathbf J}(\mathbf r',\omega)
\ee^{-\ii\kappa_0\hat{\mathbf n}\cdot\mathbf r'}\\
&=\frac1{4\pi}
(\mathbf I-\hat{\mathbf n}\hat{\mathbf n})
\widetilde{\mathbf J}(\kappa_0\hat{\mathbf n},\omega).
\end{align*}
因此
\[
|\mathbf A_{\infty,0}|^2
=\frac1{16\pi^2}
\left|(\mathbf I-\hat{\mathbf n}\hat{\mathbf n})
\widetilde{\mathbf J}\right|^2.
\]
代入\eqnrefs{eq:cl_photon_probability_far_green} 后，$1/\pi$ 与 $1/(16\pi^2)$ 合并成 $1/(16\pi^3)$，得到\eqnrefs{eq:free_space_current_radiation_probability}。
\end{derivation}

\begin{derivation}[从\eqnrefs{eq:eels_source_field_green_dp68}到\eqnrefs{eq:eels_radiative_absorptive_energy_balance}]

先由外部电流对诱导场所做的负功求总损失。时域表达式为
\begin{equation*}
 W_{\rm loss}
 =-\int\dd t\dd^3 r\,
 \mathbf J(\mathbf r,t)\cdot\mathbf E(\mathbf r,t).
\end{equation*}
利用实场的正负频共轭关系，正频谱写成
\begin{equation*}
 \frac{\dd W_{\rm loss}}{\dd\omega}
 =-\frac1\pi\operatorname{Re}
 \int\dd^3 r\,
 \widetilde{\mathbf J}^{*}(\mathbf r,\omega)
 \cdot\widetilde{\mathbf E}(\mathbf r,\omega).
\end{equation*}
代入正文\eqnrefs{eq:eels_source_field_green_dp68}，
\begin{align*}
 \frac{\dd W_{\rm loss}}{\dd\omega}
 &=-\frac1\pi\operatorname{Re}
 \left[\ii\mu_0\omega
 \langle\widetilde J|G^R|\widetilde J\rangle\right]\\
 &=\frac{\mu_0\omega}{\pi}
 \langle\widetilde J|\operatorname{Im}G^R|\widetilde J\rangle,
\end{align*}
即正文\eqnrefs{eq:eels_energy_loss_green_si}。

再计算局域介质中的不可逆吸收。正文\eqnrefs{eq:material_polarization_current_dp68} 给出
\begin{equation*}
 \widetilde{\mathbf J}_{\rm mat}
 =-\ii\omega\varepsilon_0
 (\boldsymbol\varepsilon_r-\mathbf I)\widetilde{\mathbf E}.
\end{equation*}
取正频周期平均功，
\begin{align*}
 \frac{\dd W_{\rm abs}}{\dd\omega}
 &=\frac1\pi\int_{V_{\rm mat}}\dd^3 r\,
 \operatorname{Re}
 \left[\widetilde{\mathbf J}_{\rm mat}\cdot
 \widetilde{\mathbf E}^{*}\right]\\
 &=\frac{\omega\varepsilon_0}{\pi}
 \int_{V_{\rm mat}}\dd^3 r\,
 \widetilde{\mathbf E}^{*}\cdot
 \operatorname{Im}\boldsymbol\varepsilon_r
 \cdot\widetilde{\mathbf E},
\end{align*}
得到正文\eqnrefs{eq:material_absorbed_energy_spectrum_si}。

最后证明总损失如何严格分成远场辐射与材料吸收。取一个包围全部耗散材料、外边界位于无损真空的控制体 $V$。按全文 $\ee^{-\ii\omega t}$ 约定，
\begin{equation*}
 \nabla\times\mathbf E=\ii\omega\mathbf B,
 \qquad
 \nabla\times\mathbf H=\mathbf J_{\rm ext}-\ii\omega\mathbf D,
\end{equation*}
所以共轭方程为
\begin{equation*}
 \nabla\times\mathbf H^*=\mathbf J_{\rm ext}^*+\ii\omega\mathbf D^*.
\end{equation*}
代入恒等式
\begin{equation*}
 \nabla\cdot(\mathbf E\times\mathbf H^*)
 =\mathbf H^*\cdot(\nabla\times\mathbf E)
 -\mathbf E\cdot(\nabla\times\mathbf H^*)
\end{equation*}
得到复 Poynting 方程
\begin{equation*}
 \nabla\cdot(\mathbf E\times\mathbf H^*)
 =\ii\omega\mathbf H^*\cdot\mathbf B
 -\mathbf E\cdot\mathbf J_{\rm ext}^*
 -\ii\omega\mathbf E\cdot\mathbf D^*.
\end{equation*}
介质只含电损耗并取 $\mathbf B=\mu_0\mathbf H$，第一项为纯虚数，对实部没有贡献。再写
\begin{equation*}
 \mathbf D=\varepsilon_0\boldsymbol\varepsilon_r\mathbf E,
 \qquad
 X\equiv\mathbf E^*\cdot\boldsymbol\varepsilon_r\mathbf E,
\end{equation*}
则 $\mathbf E\cdot\mathbf D^*=\varepsilon_0X^*$，从而
\begin{equation*}
 \operatorname{Re}[-\ii\omega\mathbf E\cdot\mathbf D^*]
 =-\omega\varepsilon_0\operatorname{Im}X
 =-\omega\varepsilon_0
 \mathbf E^*\cdot\operatorname{Im}\boldsymbol\varepsilon_r\cdot\mathbf E.
\end{equation*}
因此
\begin{align*}
 \operatorname{Re}\nabla\cdot(\mathbf E\times\mathbf H^*)
 ={}&-\operatorname{Re}(\mathbf J_{\rm ext}^*\cdot\mathbf E)\\
 &-\omega\varepsilon_0
 \mathbf E^*\cdot\operatorname{Im}\boldsymbol\varepsilon_r\cdot\mathbf E.
\end{align*}
对 $V$ 积分并用 Gauss 定理，
\begin{align*}
 \operatorname{Re}\oint_{\partial V}\dd\mathbf S\cdot
 (\mathbf E\times\mathbf H^*)
 ={}&-\operatorname{Re}\int_V\dd^3r\,
 \mathbf J_{\rm ext}^*\cdot\mathbf E\\
 &-\omega\varepsilon_0\int_{V_{\rm mat}}\dd^3r\,
 \mathbf E^*\cdot\operatorname{Im}\boldsymbol\varepsilon_r\cdot\mathbf E.
\end{align*}
全文的正频 Fourier 归一化给每一项共同的 $1/\pi$。于是右边第一项正是前面得到的 $\dd W_{\rm loss}/\dd\omega$，第二项是 $-\dd W_{\rm abs}/\dd\omega$；把 $\partial V$ 推到无损真空远场后，左边就是 $\dd W_{\rm rad}/\dd\omega$。因此逐频有
\begin{equation*}
 \frac{\dd W_{\rm loss}}{\dd\omega}
 =\frac{\dd W_{\rm rad}}{\dd\omega}
 +\frac{\dd W_{\rm abs}}{\dd\omega},
\end{equation*}
即正文\eqnrefs{eq:eels_radiative_absorptive_energy_balance}。这里的分解来自同一个 Maxwell 能量守恒恒等式，而不是把 EELS 与 CL 当成两套独立动力学。
\end{derivation}


